\documentclass[../zusammenfassung_AN.tex]{subfiles}
\begin{document}
	
	% -------------------------------------------------------
	\subsection{Eigenschaften von Funktionen}
	% -------------------------------------------------------
	
	\subsubsection{Symmetrieverhalten}
	
	\begin{tcolorbox}[dglbox, title={Symmetrie}]
		\textbf{Achsensymmetrie — gerade Funktion:}
		\[
		f(-x) = f(x)
		\qquad \text{Beispiele: } f(x)=\cos(x),\quad f(x)=x^2
		\]
		\textbf{Punktsymmetrie (zum Ursprung) — ungerade Funktion:}
		\[
		f(-x) = -f(x)
		\qquad \text{Beispiele: } f(x)=\sin(x),\quad f(x)=x^3
		\]
	\end{tcolorbox}
	
	\subsubsection{Monotonie}
	
	\begin{tcolorbox}[formelbox]
		Seien $x_1, x_2 \in \mathbb{D}$ mit $x_1 < x_2$. Die Funktion $f$ heisst:
		\begin{itemize}[noitemsep]
			\item \textit{monoton wachsend}, falls $f(x_1) \le f(x_2)$
			\item \textit{streng monoton wachsend}, falls $f(x_1) < f(x_2)$
			\item \textit{monoton fallend}, falls $f(x_1) \ge f(x_2)$
			\item \textit{streng monoton fallend}, falls $f(x_1) > f(x_2)$
		\end{itemize}
	\end{tcolorbox}
	
	\subsubsection{Periodizität}
	
	\begin{tcolorbox}[formelbox]
		$f$ heisst periodisch mit der Periode $p > 0$, falls für alle $x \in \mathbb{D}$:
		\[
		f(x \pm p) = f(x)
		\]
	\end{tcolorbox}
	
	\begin{itemize}[noitemsep]
		\item Mit $p$ sind auch $\pm k\cdot p$ Perioden der Funktion ($k \in \mathbb{N}_0$)
		\item Die kleinste positive Periode heisst \textbf{primitive Periode}
	\end{itemize}
	
	\subsubsection{Umkehrfunktion (Inverse Funktion)}
	
	\begin{tcolorbox}[hinweis]
		Existiert nur, wenn $f$ \textbf{streng monoton} (bijektiv) ist.
		Dann gilt: $f^{-1}(f(x)) = x$.\\
		Geometrisch: Graph von $f^{-1}$ ist die Spiegelung des Graphs von $f$ an der Geraden $y=x$.
	\end{tcolorbox}
	
	% -------------------------------------------------------
	\subsection{Ganzrationale Funktionen (Polynomfunktionen)}
	% -------------------------------------------------------
	
	\subsubsection{Lineare Funktionen (Grad 1)}
	
	\begin{tcolorbox}[formelbox]
		\renewcommand{\arraystretch}{1.5}
		\begin{tabular}{ll}
			\textbf{Normalform}         & $y = mx + b$ \\
			\textbf{Punkt-Steigungs-Form} & $y - y_1 = m(x-x_1)$ \\
			\textbf{Zwei-Punkte-Form}   & $\dfrac{y-y_1}{x-x_1} = \dfrac{y_2-y_1}{x_2-x_1}$ \\[6pt]
			\textbf{Achsenabschnittsform} & $\dfrac{x}{a}+\dfrac{y}{b}=1$ \\
		\end{tabular}
	\end{tcolorbox}
	
	\subsubsection{Quadratische Funktionen (Grad 2)}
	
	\begin{tcolorbox}[formelbox]
		\renewcommand{\arraystretch}{1.5}
		\begin{tabular}{ll}
			\textbf{Normalform}           & $y = ax^2 + bx + c$ \\
			\textbf{Scheitelpunktform}    & $y = a(x-x_S)^2 + y_S$ \\
			\textbf{Nullstellenform}      & $y = a(x-x_1)(x-x_2)$ \\
		\end{tabular}
		
		\medskip
		\textbf{Nullstellen (Lösungsformel):}
		\[
		x_{1,2} = \frac{-b \pm \sqrt{b^2-4ac}}{2a}
		\]
	\end{tcolorbox}
	
	% -------------------------------------------------------
	\subsection{Gebrochen-Rationale Funktionen}
	% -------------------------------------------------------
	
	\begin{tcolorbox}[hinweis]
		Form: $f(x) = \dfrac{p(x)}{q(x)}$, wobei $p$ und $q$ Polynome sind.
		Definitionslücken bei $q(x)=0$ (hebbar oder Pol).
	\end{tcolorbox}
	
	% -------------------------------------------------------
	\subsection{Hyperbel- und Areafunktionen}
	% -------------------------------------------------------
	
	\begin{tcolorbox}[dglbox, title={Hyperbolische Funktionen}]
		\[
		\sinh(x) = \frac{e^x - e^{-x}}{2} \qquad
		\cosh(x) = \frac{e^x + e^{-x}}{2} \qquad
		\tanh(x) = \frac{\sinh(x)}{\cosh(x)} = \frac{e^x - e^{-x}}{e^x + e^{-x}}
		\]
		\textbf{Hyperbolischer Pythagoras:}
		\[
		\cosh^2(x) - \sinh^2(x) = 1
		\]
	\end{tcolorbox}
	
\end{document}